\documentclass[11pt]{article}
\usepackage[margin=1in]{geometry}
\usepackage{amsmath,amssymb,amsthm,mathtools}
\usepackage{graphicx}
\usepackage{xcolor}
\usepackage{microtype}
\usepackage[colorlinks=true,linkcolor=teal!70!black,citecolor=teal!70!black,urlcolor=teal!70!black]{hyperref}

\newtheorem{theorem}{Theorem}
\newtheorem{corollary}[theorem]{Corollary}
\newcommand{\cF}{\mathcal F}

\title{Wolff Reiteration Beyond Lattices:\\
A Compatible Complemented-Dilation Theorem}
\author{Run \texttt{fa\_banach\_001}, \texttt{agent\_lane\_14}}
\date{12 August 2026}

\begin{document}
\maketitle

\section*{Status and source question}

\textbf{Status:} candidate substantial partial result, likely valid.
Egert and Kosmala~\cite{EgertKosmala2024} prove a Wolff-reiteration theorem
for $A$-convex quasi-Banach function spaces.  On source PDF page~2,
Remark~1.3 asks for the same conclusion for more general $A$-convex
quasi-Banach spaces.

\begin{center}
\includegraphics[width=0.97\textwidth]{figures/open_problem_crop.png}
\end{center}

The theorem below gives the full Wolff conclusion for a class of spaces which
need not themselves be function lattices: compatible common projection ranges
of a function-space quadruple, provided the complementary quadruple obeys the
same two local interpolation identities.  The condition is automatic for a
stationary complement.

Fix parameters
\begin{equation}\label{eq:parameters}
 0<\theta<\eta<1,\qquad \theta=\lambda\eta,
 \qquad \eta=(1-\mu)\theta+\mu.
\end{equation}
All complex brackets below denote the outer complex method used in the source.

\section*{Proof intuition}

A common bounded projection splits each ambient function space as
$Y_j=X_j\oplus W_j$.  Complex interpolation respects such finite direct sums.
Thus, if both the desired range spaces $X_j$ and their complements $W_j$ obey
the two local identities, those identities recombine into the hypotheses of
Egert--Kosmala's function-space theorem for the $Y_j$.  That theorem supplies
the two global identities upstairs.  Projecting back extracts exactly the
desired identities for the possibly non-lattice spaces $X_j$.

The complement hypothesis is not bookkeeping: without information about the
$W_j$, the two local identities for the $X_j$ do not lift to the $Y_j$.  This
pinpoints what the projection argument can and cannot prove.

\section*{Complemented-dilation theorem}

\begin{theorem}[Compatible complemented-lattice dilation]
\label{thm:dilation}
Let $Y_0,Y_1,Y_2,Y_3$ be $A$-convex quasi-Banach function spaces over one
separable, $\sigma$-finite measure space.  Suppose that either $Y_0$ or $Y_3$
is separable and that $Y_0\cap Y_3$ is dense in $Y_1$ and $Y_2$.

Let $P$ be a linear projection on a common ambient topological vector space
containing the $Y_j$, and suppose that $P|_{Y_j}$ is bounded for
$j=0,1,2,3$.  Put
\[
 X_j=P(Y_j),\qquad W_j=(I-P)(Y_j),\qquad j=0,1,2,3,
\]
with the inherited quotient (equivalently, subspace) quasi-norms.  Assume,
compatibly in the common ambient spaces, that
\begin{align}
 X_1&=[X_0,X_2]_\lambda,& X_2&=[X_1,X_3]_\mu,
 \label{eq:xlocal}\\
 W_1&=[W_0,W_2]_\lambda,& W_2&=[W_1,W_3]_\mu.
 \label{eq:wlocal}
\end{align}
Then
\begin{equation}\label{eq:conclusion}
 X_1=[X_0,X_3]_\theta,\qquad X_2=[X_0,X_3]_\eta
\end{equation}
up to equivalence of quasi-norms.
\end{theorem}

\begin{proof}
Since $P$ and $I-P$ are bounded on $Y_j$, the map
\[
 Y_j\longrightarrow X_j\oplus W_j,\qquad y\longmapsto (Py,(I-P)y)
\]
is an isomorphism with inverse $(x,w)\mapsto x+w$.  These maps are compatible
for all four indices, so we may identify
\begin{equation}\label{eq:splitY}
 Y_j=X_j\oplus W_j
\end{equation}
topologically.

We first record the elementary direct-sum identity
\begin{equation}\label{eq:directsum}
 [A_0\oplus B_0,A_1\oplus B_1]_s
   =[A_0,A_1]_s\oplus[B_0,B_1]_s,\qquad 0<s<1.
\end{equation}
Indeed, an admissible function $F$ for the couple on the left splits
coordinatewise as $F=(F_A,F_B)$, where $F_A$ and $F_B$ are admissible for the
two component couples.  Conversely, two componentwise admissible functions
combine into one.  The boundary and ambient-sum quasi-norms are equivalent to
the maximum (or any fixed finite direct-sum quasi-norm) of their component
quasi-norms.  Taking infima over functions with prescribed value proves
\eqref{eq:directsum} with equivalence of quasi-norms.  This is also the usual
retract principle for the complex method; compare
Kalton--Mayboroda--Mitrea~\cite[Lemma~7.11]{KMM2007}.

Using \eqref{eq:splitY}, \eqref{eq:directsum}, and the two identities at
parameter $\lambda$, we obtain
\[
 [Y_0,Y_2]_\lambda
  =[X_0,X_2]_\lambda\oplus[W_0,W_2]_\lambda
  =X_1\oplus W_1=Y_1.
\]
The identities at parameter $\mu$ similarly give
\[
 [Y_1,Y_3]_\mu=Y_2.
\]
All hypotheses of Egert--Kosmala's Theorem~1.2 now hold for the function-space
quadruple $(Y_0,Y_1,Y_2,Y_3)$.  Hence
\begin{equation}\label{eq:upstairs}
 Y_1=[Y_0,Y_3]_\theta,\qquad Y_2=[Y_0,Y_3]_\eta.
\end{equation}

Finally, \eqref{eq:directsum} applied to $Y_0=X_0\oplus W_0$ and
$Y_3=X_3\oplus W_3$ shows, for every $0<s<1$, that
\[
 P[Y_0,Y_3]_s=[X_0,X_3]_s.
\]
Applying $P$ to the two identities in \eqref{eq:upstairs}, and using
$P(Y_1)=X_1$ and $P(Y_2)=X_2$, proves \eqref{eq:conclusion}.
\end{proof}

\begin{corollary}[Stationary complement]
Under the hypotheses of Theorem~\ref{thm:dilation}, suppose that there is one
quasi-Banach space $W$ such that the four $W_j$ are compatibly equal to $W$
with equivalent quasi-norms.  Then the local complement identities
\eqref{eq:wlocal} are automatic, so the conclusion
\eqref{eq:conclusion} follows from \eqref{eq:xlocal} alone.
\end{corollary}

\begin{proof}
For every $0<s<1$, $[W,W]_s=W$ with equivalent quasi-norms.  Constant
admissible functions give $W\hookrightarrow[W,W]_s$, while the definition of
admissibility and the maximum principle for the $A$-convex ambient space give
the reverse inclusion.  Substitution yields \eqref{eq:wlocal}, and
Theorem~\ref{thm:dilation} applies.
\end{proof}

\section*{Verification and scope}

The proof is formal but not vacuous.  A bounded projection range of a
quasi-Banach function lattice need not be an order ideal and need not carry a
compatible lattice structure, so the theorem reaches abstract $A$-convex
spaces outside the literal hypotheses of the source theorem.  The
$A$-convexity of each $X_j$ follows independently because it is isomorphic to
a closed subspace of the $A$-convex space $Y_j$.  Moreover,
$P(Y_0\cap Y_3)\subseteq X_0\cap X_3$ is dense in $X_1$ and $X_2$, so the
density condition expected in the source question is inherited.

Each use of complex interpolation in the proof is justified directly from
admissible functions.  In particular, no vector-valued integration, duality,
or unproved ordinary quasi-Banach reiteration theorem is being smuggled into
the argument.

The result does not settle arbitrary $A$-convex quadruples.  Eight focused
upgrade attempts were made: gluing two local scales, applying the
Kalton--Mitrea stability theorem, Banach-envelope/convexification reduction,
constructing a non-$A$-convex-sum counterexample, passing to the inner method,
common unconditional coordinates, complemented dilation, and finally removal
of the complement hypothesis.  The first six routes meet the same obstruction:
for general quasi-Banach spaces the outer method need not be known to equal the
core-generated inner method, and the Banach proof of reverse reiteration uses
integration or duality unavailable here.  The last route shows that a common
projection alone supplies no identities among $W_0,W_1,W_2,W_3$; hence the
complement condition cannot be deleted by this argument.

Bounded exact-title, arXiv-id, author, and keyword searches through
12 August 2026 used the phrases \emph{Wolff-reiteration},
\emph{$A$-convex quasi-Banach}, \emph{inner complex interpolation},
\emph{outer complex interpolation}, and \emph{reiteration theorem}.  They
found no later general proof or counterexample.  The theorem above appears to
be a direct new corollary of known results rather than a new interpolation
principle; novelty confidence is modest.

\section*{Human-review recommendation}

Check that the four copies of $P$ really are restrictions of one compatible
projection in the intended application, rather than merely unrelated
isomorphic projections.  Then verify the two complement identities, or the
stationary-complement condition.  Those are the only application-specific
points; the direct-sum and projection proof is otherwise exact.  A full upgrade
would require either a general reverse-reiteration theorem for the outer
quasi-Banach complex method or a counterexample showing it fails.

\begin{thebibliography}{9}
\bibitem{EgertKosmala2024}
M.~Egert and B.~W. Kosmala,
\emph{A Note on Complex Interpolation of Quasi-Banach Function Spaces},
arXiv:2412.12769 (2024); published version, Studia Math. (2026).

\bibitem{KMM2007}
N.~J. Kalton, S.~Mayboroda, and M.~Mitrea,
\emph{Interpolation of Hardy--Sobolev--Besov--Triebel--Lizorkin spaces and
applications to problems in partial differential equations},
in \emph{Interpolation Theory and Applications}, Contemp. Math. 445,
Amer. Math. Soc., 2007, 121--177.

\bibitem{KaltonMitrea1998}
N.~J. Kalton and M.~Mitrea,
\emph{Stability results on interpolation scales of quasi-Banach spaces and
applications}, Trans. Amer. Math. Soc. 350 (1998), 3903--3922.
\end{thebibliography}

\end{document}
